
\chapter{Noga Alon (2024)}
\index{Alon, Noga}

\begin{quote}
\it ``Noga Alon has been one of the most influential combinatorialists of the modern era. His work, combining deep mathematical insight with algorithmic thinking, has created entire new areas of research at the intersection of combinatorics, computer science, and number theory. The combinatorial Nullstellensatz, the algorithmic Lov\'asz Local Lemma, and his contributions to graph theory and additive combinatorics have become essential tools across mathematics.'' --- Wolf Prize Citation, 2024
\end{quote}

\section{Early Life and Career}

Noga Alon was born on April 29, 1956, in Haifa, Israel. He studied at the Hebrew University of Jerusalem, receiving his Ph.D.\ in 1983 under the supervision of Vitali Milman. He has held positions at Tel Aviv University and the Institute for Advanced Study at Princeton, and is currently a professor at the Hebrew University of Jerusalem and New York University. He was awarded the Wolf Prize in Mathematics in 2024.

\section{Main Contributions}

Noga Alon (born 1956) is an Israeli mathematician whose work has transformed combinatorics and its connections to theoretical computer science. He was awarded the Wolf Prize in Mathematics in 2024 alongside Adi Shamir.

The Wolf Prize citation reads: ``for their pioneering contributions to the mathematical foundations of computer science, and for their profound and lasting impact on combinatorics and number theory.''

Alon's core contributions include:
\begin{itemize}
    \item \textbf{Combinatorial Nullstellensatz\index{combinatorial Nullstellensatz}:} A powerful algebraic tool for proving the existence of combinatorial configurations by showing that certain polynomials cannot vanish at all points of a product set.
    \item \textbf{Graph Theory and Colorings:} The Alon--Tarsi theorem on the chromatic number of graphs and the Alon--Boppana bound on the eigenvalues of regular graphs.
    \item \textbf{The Lov\'asz Local Lemma:} The development of algorithmic versions of the Lov\'asz Local Lemma, making it a practical tool for constructing large combinatorial objects.
    \item \textbf{Additive Combinatorics:} The Alon--Bourgain--Conlon--Gowers--Ruzsa theory of sum-free sets and the Erdos--Ginzburg--Ziv theorem (zero-sum combinatorics).
    \item \textbf{Probabilistic Method:} Foundational contributions to the probabilistic method in combinatorics, including the Alon--Janson--Rucinski inequality for concentration of random variables.
    \item \textbf{Pseudorandomness and Expanders:} The construction of optimal expander graphs\index{expander graph} and the theory of pseudorandom graphs.
    \item \textbf{Combinatorial Number Theory:} The Alon--F\"uredi theorem on covering numbers and results on the sunflower lemma.
\end{itemize}

\section{Origin of the Problem}

\subsection{The Combinatorial Nullstellensatz}

The polynomial method in combinatorics uses algebraic properties of polynomials to prove combinatorial results. A fundamental question: given a polynomial $P(x_1, \ldots, x_n)$ and a product set $S_1 \times \cdots \times S_n \subset \F^n$, if $P$ vanishes on all points of the product set, what does this imply about the coefficients of $P$?

The classical Nullstellensatz (Hilbert) is about zeros of polynomials in algebraically closed fields. The combinatorial Nullstellensatz provides a simple but powerful criterion for non-vanishing on product sets.

\subsection{The Lov\'asz Local Lemma}

The Lov\'asz Local Lemma (LLL) is a probabilistic tool for proving that a set of unlikely events can all be avoided simultaneously, provided each event depends on few others. The original lemma was non-constructive. The question: can one find such configurations algorithmically?

\subsection{Graph Eigenvalues and Expanders}

The eigenvalues of the adjacency matrix of a graph encode much of its combinatorial structure. A fundamental bound (Alon--Boppana) states that any infinite family of $d$-regular graphs has an eigenvalue at least $2\sqrt{d-1}$.

\section{How It Was Solved and the Intuitions/Tools Needed}

\subsection{The Combinatorial Nullstellensatz}

The theorem states: let $\F$ be a field and let $P \in \F[x_1, \ldots, x_n]$. Suppose $\deg P = \sum_{i=1}^n t_i$ and the coefficient of $\prod x_i^{t_i}$ is nonzero. Then for any subsets $S_i \subset \F$ with $|S_i| > t_i$, there exists $(s_1, \ldots, s_n) \in \prod S_i$ such that $P(s_1, \ldots, s_n) \neq 0$.

The proof uses the fact that monomials are linearly independent functions on the product set.

\subsection{Algorithmic Lov\'asz Local Lemma}

Alon developed a general method to make the LLL constructive:
\begin{enumerate}
    \item Use the LLL to prove that a random configuration avoids all bad events with positive probability.
    \item Show that the expectation-computation technique (using the LLL condition) can be turned into an efficient algorithm by resampling variables.
    \item The Moser--Tardos algorithm (building on Alon's ideas) provides a simple, efficient algorithm.
\end{enumerate}

\section{Mathematical Theory}

\subsection{The Combinatorial Nullstellensatz}

\begin{theorem}[Combinatorial Nullstellensatz, Alon 1999]
Let $\F$ be a field and $P \in \F[x_1, \ldots, x_n]$. Let $S_1, \ldots, S_n$ be finite subsets of $\F$ with $|S_i| > \deg_{x_i}(P)$. If $P(s_1, \ldots, s_n) = 0$ for all $(s_1, \ldots, s_n) \in \prod S_i$, then $P \equiv 0$.
\end{theorem}

\begin{corollary}[Nonvanishing Criterion]
If $\deg P = \sum t_i$ and the coefficient of $\prod x_i^{t_i}$ in $P$ is nonzero, then for any $S_i$ with $|S_i| > t_i$, there exists $s \in \prod S_i$ with $P(s) \neq 0$.
\end{corollary}

Applications include:
\begin{itemize}
    \item \textbf{Graph Coloring:} Alon--Tarsi theorem on list colorings.
    \item \textbf{Additive Combinatorics:} Zero-sum problems.
    \item \textbf{Number Theory:} Existence of solutions to polynomial equations.
\end{itemize}

\subsubsection*{The Alon--F\"uredi Theorem}

\begin{theorem}[Alon--F\"uredi, 1993]
Let $\F$ be a field and let $S_1, \ldots, S_n$ be finite subsets of $\F$ with $|S_i| = s_i$. Let $P \in \F[x_1, \ldots, x_n]$ be a nonzero polynomial of the form
\[
P(x_1, \ldots, x_n) = c_{\alpha} \prod_{i=1}^n x_i^{\alpha_i} + \text{(terms of lower total degree in the lexicographic ordering)},
\]
where $c_\alpha \neq 0$ and $\deg P = \sum_{i=1}^n \alpha_i$. Then the number of nonzeros of $P$ on $S_1 \times \cdots \times S_n$ is at least
\[
\prod_{i=1}^n (s_i - \alpha_i).
\]
In particular, if $\deg_{x_i}(P) < |S_i|$ for all $i$, then $P$ has at least $\prod_{i=1}^n (|S_i| - \deg_{x_i}(P))$ nonzeros on the product set.
\end{theorem}

\begin{remark}
The Alon--F\"uredi theorem strengthens the basic combinatorial Nullstellensatz by providing a \emph{quantitative lower bound} on the number of nonzeros, rather than merely asserting their existence. When $|S_i| = s$ and each $\deg_{x_i}(P) \leq t$, the bound becomes $(s - t)^n$, which is often sharp.
\end{remark}

\subsubsection*{Transversals in Latin Squares}

A \emph{Latin square} of order $n$ is an $n \times n$ array filled with symbols from $[n] = \{1, \ldots, n\}$ such that each symbol appears exactly once in each row and each column. A \emph{transversal} is a set of $n$ cells, one in each row and one in each column, containing all $n$ symbols.

\begin{theorem}[Alon, 1999]
Every Latin square of odd order has a transversal.
\end{theorem}

\begin{proof}
Let $L = (a_{ij})$ be an $n \times n$ Latin square with $n$ odd, where $a_{ij} \in [n]$. Consider the polynomial over $\F_n$ (where $n$ is prime for simplicity; the general case uses a reduction):
\[
P(x_1, \ldots, x_n) = \prod_{i=1}^{n} \prod_{1 \leq j < k \leq n} (x_j - x_k) \cdot \prod_{c=1}^{n} \left(\sum_{i=1}^{n} a_{i,x_i} - c\right).
\]
By the combinatorial Nullstellensatz, one shows that the coefficient of the top monomial $x_1^{n-1} \cdots x_n^{n-1}$ in a carefully constructed auxiliary polynomial is nonzero. The oddness of $n$ ensures that this coefficient survives the algebraic manipulations, yielding a nonzero evaluation point, which corresponds to a transversal.
\end{proof}

\begin{remark}
The existence of transversals in Latin squares of even order remains a major open problem, closely related to Ryser's conjecture. The polynomial method introduced by Alon provides the strongest known partial results.
\end{remark}

\subsection{The Alon--Boppana Bound}

\begin{theorem}[Alon--Boppana, 1986]
Let $G$ be a $d$-regular graph on $n$ vertices with adjacency eigenvalues $\lambda_1 = d \geq \lambda_2 \geq \cdots \geq \lambda_n$. Then:
\[
\lambda_2 \geq 2\sqrt{d-1} - o_n(1)
\]
where $o_n(1) \to 0$ as $n \to \infty$.
\end{theorem}

This bound shows that the spectral gap of a $d$-regular graph is at most $d - 2\sqrt{d-1}$. Ramanujan graphs achieve this bound.

\subsection{Expander Graphs and their Applications}

Expander graphs are sparse graphs with exceptionally strong connectivity properties. They arise naturally in theoretical computer science, number theory, and combinatorics, and Alon's work was instrumental in developing their theory.

\begin{definition}[Edge Expansion]
Let $G = (V, E)$ be a $d$-regular graph on $n$ vertices. The \textbf{edge expansion} (or \textbf{Cheeger constant}) of $G$ is
\[
h(G) = \min_{\substack{S \subset V \\ 0 < |S| \leq n/2}} \frac{|\partial S|}{|S|},
\]
where $\partial S$ denotes the set of edges with exactly one endpoint in $S$. A family of $d$-regular graphs $\{G_n\}$ is an \textbf{expander family} if $h(G_n) \geq \epsilon > 0$ for some constant $\epsilon$ independent of $n$.
\end{definition}

\begin{definition}[Spectral Gap]
For a $d$-regular graph $G$ with adjacency eigenvalues $d = \lambda_1 \geq \lambda_2 \geq \cdots \geq \lambda_n$, the \textbf{spectral gap} is $d - \lambda_2$. The graph is an expander if and only if the spectral gap is bounded away from zero: $d - \lambda_2 \geq \delta > 0$ for some constant $\delta$, uniformly over the family.
\end{definition}

The connection between spectral gap and edge expansion is given by the Cheeger inequality:
\[
\frac{d - \lambda_2}{2} \leq h(G) \leq \sqrt{2d(d - \lambda_2)}.
\]
Thus the spectral characterization and the combinatorial definition of expanders are equivalent up to constant factors.

\begin{theorem}[Expander Mixing Lemma]
Let $G$ be an $(n, d, \lambda)$-graph. For any two subsets $A, B \subseteq V(G)$, the number of edges between $A$ and $B$ satisfies
\[
\left| e(A, B) - \frac{d \cdot |A| \cdot |B|}{n} \right| \leq \lambda \sqrt{|A| \cdot |B|},
\]
where $e(A,B)$ denotes the number of edges with one endpoint in $A$ and the other in $B$.
\end{theorem}

This lemma makes precise the sense in which expanders ``look random'' locally.

\subsubsection*{Ramanujan Graphs}

The Alon--Boppana bound shows that any infinite family of $d$-regular graphs must have $\lambda_2 \geq 2\sqrt{d-1} - o(1)$. A $d$-regular graph with all nontrivial eigenvalues satisfying $|\lambda_i| \leq 2\sqrt{d-1}$ is called a \textbf{Ramanujan graph}. Such graphs are optimal expanders.

\begin{theorem}[Lubotzky--Phillips--Sarnak, 1988; Margulis, 1988]
For every prime power $q \equiv 1 \pmod{4}$, there exists a family of $(q+1)$-regular Ramanujan graphs on $q^2(q^2+1)/2$ or $q^2(q^2-1)/2$ vertices. The construction uses Cayley graphs of $\PSL(2, \F_q)$ or $\PGL(2, \F_q)$ with generating sets chosen via the Hilbert symbol.
\end{theorem}

The proof constructs the graph as a Cayley graph $\Cay(\Gamma, S)$ where $\Gamma = \PSL(2, \F_q)$ and $S$ is a carefully chosen symmetric set of generators. The eigenvalues are then bounded using the representation theory of $\PSL(2, \F_q)$, specifically the Weil representation and bounds on Kloosterman sums. The Ramanujan conjecture for cusp forms (proved by Deligne) is used to control the character sums that arise.

\begin{remark}
The existence of Ramanujan graphs was a major motivation for the study of algebraic graph theory and arithmetic expanders. Gaber and Galil (1981) gave explicit algebraic constructions achieving constant spectral gap, but the Ramanujan graphs of Lubotzky--Phillips--Sarnak achieved the information-theoretically optimal bound.
\end{remark}

\subsubsection*{Applications to Error-Correcting Codes}

Expander graphs play a central role in the construction of efficient error-correcting codes:

\begin{theorem}[Expander Codes, Sipser--Spielman, 1996]
There exist families of linear codes constructed from expander graphs that achieve constant rate, constant relative distance, and can be decoded in linear time using a simple belief-propagation algorithm.
\end{theorem}

The key idea is to view a bipartite expander graph as the parity-check matrix of a code. The expansion property guarantees that the minimum distance is large (at least a constant fraction of the block length), while the sparsity of the graph enables efficient decoding. Specifically, if $G$ is a bipartite $(d, c)$-regular expander with spectral gap $\lambda$, then the corresponding code has relative distance at least $1 - c/d - \lambda/d$.

Expander-based codes include:
\begin{itemize}
    \item \textbf{LDPC codes} (Low-Density Parity-Check): codes whose parity-check matrices are sparse bipartite graphs. Expander-based LDPC codes achieve capacity on binary symmetric channels.
    \item \textbf{Reed--Muller codes on expanders}: using expanders to derandomize the construction of good codes.
    \item \textbf{Algebraic-geometric codes}: expanders appear in the decoding of AG codes via expander-based list decoding.
\end{itemize}

\subsubsection*{Applications to Derandomization}

Expander graphs are a cornerstone of derandomization:

\begin{theorem}[Random Walks on Expanders]
Let $G$ be a $d$-regular expander on $n$ vertices with second eigenvalue $\lambda_2$. After $O(\log n)$ steps, a random walk on $G$ produces a distribution whose statistical distance from the uniform distribution is at most $n^{-c}$ for any desired constant $c > 0$.
\end{theorem}

This mixing property is used to construct:
\begin{itemize}
    \item \textbf{Pseudorandom generators}: given a hard function, construct a generator that fools any small circuit. Expanders yield generators with short seed length.
    \item \textbf{Extractor functions}: functions that extract nearly uniform randomness from weak random sources. Leftover hash lemma combined with expander walks yields near-optimal extractors.
    \item \textbf{Samplers}: procedures that estimate the fraction of satisfying inputs to a circuit using few queries. Expander-based samplers achieve $O(1/\epsilon^2)$ samples for $\epsilon$-approximation.
\end{itemize}

\subsection{The Alon--Tarsi Theorem}

\begin{theorem}[Alon--Tarsi, 1992]
Let $G$ be a graph and let $k$ be an integer. For each edge, assign an orientation and a color in $\{1, \ldots, k\}$. If the number of Eulerian orientations is odd, then $G$ is $k$-choosable.
\end{theorem}

This theorem connects parity of Eulerian orientations to list coloring.

\subsection{Additive Combinatorics}

\begin{theorem}[Erdos--Ginzburg--Ziv Theorem, Alon generalization]
For any $2n-1$ integers, there exists a subset of size $n$ whose sum is divisible by $n$. Alon extended this to abelian groups and provided algorithmic proofs.
\end{theorem}

\begin{theorem}[Alon--F\"uredi, 1993]
The covering number $C(n,k)$ (minimum number of $k$-dimensional subspaces needed to cover $\F_q^n$) satisfies:
\[
C(n,k) \geq \frac{q^n - 1}{q^k - 1}
\]
with equality only for the trivial covering by all $\frac{q^n-1}{q^k-1}$ subspaces.
\end{theorem}

\subsection{Pseudorandom Graphs}

\begin{definition}[$(n,d,\lambda)$-Graph]
An \textbf{$(n,d,\lambda)$-graph} is a $d$-regular graph on $n$ vertices with all eigenvalues (except $d$) bounded by $\lambda$ in absolute value.
\end{definition}

\begin{theorem}[Alon, 1986]
An $(n,d,\lambda)$-graph has strong pseudorandom properties: for any subset $S$ of vertices, the number of edges inside $S$ is approximately $d|S|^2/(2n)$, with error controlled by $\lambda$.
\end{theorem}

\subsection{Pseudorandomness and Derandomization}

The method of conditional expectations and the use of expanders provide powerful tools for removing randomness from probabilistic algorithms.

\subsubsection*{The Method of Conditional Expectations}

The method of conditional expectations is a derandomization technique that converts a probabilistic existence proof into a deterministic algorithm.

\begin{theorem}[Method of Conditional Expectations]
Let $X$ be a nonneg random variable and suppose $\E[X] \geq 1$. Then there exists a deterministic procedure that, given oracle access to the conditional expectations $\E[X \mid x_1, \ldots, x_k]$, finds a sequence $(x_1, \ldots, x_k)$ with $X(x_1, \ldots, x_k) \geq 1$.
\end{theorem}

\begin{proof}
We sequentially choose $x_1, \ldots, x_k$. Having chosen $x_1, \ldots, x_{i-1}$, we select $x_i$ to be any value in the domain such that
\[
\E[X \mid x_1, \ldots, x_{i-1}, x_i] \geq \E[X \mid x_1, \ldots, x_{i-1}].
\]
Such a value exists by the law of total expectation. Since $\E[X] \geq 1$, the final conditional expectation satisfies $\E[X \mid x_1, \ldots, x_k] \geq 1$, and since $X$ is determined by $(x_1, \ldots, x_k)$, we have $X(x_1, \ldots, x_k) \geq 1$.
\end{proof}

\begin{remark}
The computational bottleneck is evaluating the conditional expectations, which may be expensive. For example, derandomizing the algorithmic Lov\'asz Local Lemma requires computing conditional probabilities of complex events. Expander-based techniques can reduce the number of random bits needed, thereby reducing the number of conditional expectation evaluations.
\end{remark}

\subsubsection*{Derandomized Algorithms via Expanders}

Expander graphs enable the replacement of truly random bits with pseudorandom bits derived from expander walks:

\begin{theorem}[Expander Walk Lemma]
Let $G$ be a $d$-regular expander on $n$ vertices with second largest eigenvalue $\lambda_2 \leq \lambda$. Let $f : V \to \{0, 1\}$ be any function with $\mu = \E_{v}[f(v)]$ under the uniform distribution. For a random walk $v_0, v_1, \ldots, v_k$ of length $k$ on $G$:
\[
\Pr\left[\left|\frac{1}{k}\sum_{i=1}^{k} f(v_i) - \mu\right| > t\right] \leq 2\exp\left(-\frac{k(d - \lambda)^2 t^2}{2d^2}\right).
\]
\end{theorem}

This concentration bound is nearly as strong as the Chernoff bound for independent random variables, but requires only $\log_d n$ random bits per step (to specify the edge chosen at each step), compared to $\log n$ bits for independent sampling from $[n]$. The total number of random bits is $k \log d$, independent of $n$.

The applications include:
\begin{itemize}
    \item \textbf{Optimal extractors}: Buildhash--Zuckerman (2006) constructed near-optimal randomness extractors using expander walks, achieving seed length $O(\log n) + O(\log(1/\epsilon))$ for output length $m$ and statistical distance $\epsilon$.
    \item \textbf{Space-efficient sampling}: Nisan--Zuckerman (1995) used expander walks to construct samplers that use $O(\log n)$ bits of space, compared to $O(n)$ for naive approaches.
    \item \textbf{Derandomized MAX-SAT}: Expander-based conditioning gives $O(1)$-approximation algorithms for MAX-SAT using $O(\log n)$ random bits.
\end{itemize}

\subsubsection*{The Zig-Zag Product of Graphs}

The zig-zag product, introduced by Reingold-- Vadhan--Wigderson (2000), provides an operation on graphs that preserves regularity and spectral gap while reducing the number of vertices. It is the key ingredient in the construction of deterministic logspace algorithms for undirected connectivity (SL = L).

\begin{definition}[Zig-Zag Product]
Let $G$ be a $D$-regular graph on $n$ vertices and let $H$ be a $d$-regular graph on $D$ vertices. The \textbf{zig-zag product} $G \circ H$ is a $d^2$-regular graph on $n \cdot d$ vertices defined as follows. Replace each vertex $v$ of $G$ by a copy $H_v$ of $H$. An edge $(u, v)$ in $G$ is replaced by a bijection between the vertices of $H_u$ and $H_v$. A step in $G \circ H$ from a vertex $i$ in $H_u$ consists of: (1) follow the edge within $H_u$ determined by the current state (the ``zig''); (2) cross to $H_v$ via the edge $(u,v)$ in $G$ if $i$ is in the appropriate partition class; (3) follow an edge within $H_v$ (the ``zag'').
\end{definition}

\begin{theorem}[Reingold--Vadhan--Wigderson, 2000]
The zig-zag product satisfies:
\begin{enumerate}
    \item If $G$ is $D$-regular and $H$ is $d$-regular, then $G \circ H$ is $d^2$-regular.
    \item $\lambda(G \circ H) \leq \lambda(G) + \lambda(H)$.
    \item If $G_n$ is a family of $D$-regular expanders with $D = O(1)$ and $\lambda(G_n) \leq 1/2$, then iterating the zig-zag product with a fixed small graph $H$ yields a family of expanders with $n$ growing exponentially at each step.
\end{enumerate}
\end{theorem}

The zig-zag product allows the iterative construction of expanders from a single small constant-sized expander, yielding the graph $\tilde{O}(n)$ that is an expander on $n$ vertices with $n$ growing as $\Delta^{2^t}$ after $t$ iterations. This construction is the foundation of Reingold's deterministic logspace algorithm for $s$-$t$ connectivity in undirected graphs.

\section{Impact on Science and Technology}

\subsection{In Mathematics}

\begin{enumerate}
    \item \textbf{Combinatorics:} The combinatorial Nullstellensatz, Alon--Tarsi theorem, and the probabilistic method are standard tools.
    \item \textbf{Additive Combinatorics:} Foundations for sum-free sets, zero-sum problems.
    \item \textbf{Number Theory:} Applications of the polynomial method to number-theoretic problems.
\end{enumerate}

\subsection{In Computer Science}

\begin{enumerate}
    \item \textbf{Algorithms:} The algorithmic Lov\'asz Local Lemma is used in constraint satisfaction, graph coloring, and scheduling.
    \item \textbf{Complexity Theory:} Expander graphs and pseudorandomness are essential for derandomization.
    \item \textbf{Coding Theory:} Expanders are used in error-correcting codes (expander codes).
    \item \textbf{Network Design:} Expanders provide optimal connectivity for communication networks.
\end{enumerate}

\section{Frequently Asked Questions}

\subsection*{Q1: What is the combinatorial Nullstellensatz?}
A theorem stating that if a polynomial $P$ with a given leading monomial vanishes on a product set of size larger than the degree, then $P$ must be identically zero. It is used to prove existence of combinatorial configurations.

\subsection*{Q2: What is the Alon--Boppana bound?}
A lower bound on the second eigenvalue of a $d$-regular graph: $\lambda_2 \geq 2\sqrt{d-1} - o(1)$. It shows that Ramanujan graphs have optimal spectral gap.

\subsection*{Q3: What is the algorithmic Lov\'asz Local Lemma?}
A constructive version of the LLL that provides efficient algorithms to find configurations avoiding bad events, not just prove their existence.

\subsection*{Q4: What is an expander graph?}
A sparse graph with strong connectivity properties. Expanders are used in network design, coding theory, and derandomization.

\subsection*{Q5: What should an undergraduate study to understand Alon's work?}
Combinatorics (graph theory, extremal combinatorics), probability theory, linear algebra. Recommended: \textit{The Probabilistic Method} (Alon--Spencer), \textit{Additive Combinatorics} (Tao--Vu), and \textit{Graph Theory} (Diestel).

\section{Related Chapters}
See also Chapter~34 (Lovász) on combinatorics and graph theory and Chapter~68 (Shamir) on theoretical computer science.

\section*{Exercises for the Reader}
\addcontentsline{toc}{section}{Exercises}

\begin{enumerate}
    \item [I] Use the combinatorial Nullstellensatz to prove that any graph of maximum degree $\Delta$ is $(\Delta+1)$-colorable.
    \item [A] Prove the Alon--Boppana bound for $d$-regular trees.
    \item [I] Show that the $(n,d,\lambda)$-graph property implies the expander mixing lemma.
    \item [I] Apply the probabilistic method to show existence of a $k$-coloring of a graph with high girth.
    \item [I] Prove the Erdos--Ginzburg--Ziv theorem for $n$ prime.
    \item [B] Show that the Alon--Tarsi theorem implies that even cycles are 2-choosable.
    \item [A] Construct a Ramanujan graph using Cayley graphs of $\PSL(2,q)$.
\end{enumerate}

\section*{Timeline}
\addcontentsline{toc}{section}{Timeline}

\begin{tabularx}{\linewidth}{|p{1.5cm}|>{\raggedright\arraybackslash}X|}
\hline
\textbf{Year} & \textbf{Event} \\ \hline
1956 & Born in Haifa, Israel \\ \hline
1975--79 & B.Sc., M.Sc., Ph.D.\ from Hebrew University \\ \hline
1985 | Alon--Boppana bound \\ \hline
1992 | Alon--Tarsi theorem \\ \hline
1995 | Moves to Tel Aviv University \\ \hline
1999 | Combinatorial Nullstellensatz \\ \hline
2000 | Alon--Spencer \textit{The Probabilistic Method} (1st ed.) \\ \hline
2008 | Israel Prize in Mathematics \\ \hline
2010 | Chairman of the Israeli Council for Higher Education \\ \hline
2021 | Moves to Princeton University \\ \hline
2024 | Wolf Prize in Mathematics \\ \hline
\end{tabularx}

\section*{Glossary}
\addcontentsline{toc}{section}{Glossary}

\begin{description}
    \item[Additive combinatorics] The study of additive structures in sets of integers or group elements.
    \item[Combinatorial Nullstellensatz] An algebraic tool for proving non-vanishing of polynomials on product sets.
    \item[Expander graph] A graph with high connectivity relative to its sparsity.
    \item[Lov\'asz Local Lemma] A probabilistic lemma for avoiding unlikely collections of bad events.
    \item[Probabilistic method] A method for proving existence of combinatorial objects using probability.
    \item[Pseudorandom graph] A graph that mimics the properties of a random graph.
    \item[Spectral gap] The difference between the largest and second-largest eigenvalues of a graph.
    \item[Zero-sum] The study of subsets summing to zero in abelian groups.
\end{description}

\section{Citations}\subsection*{Primary Sources}
\begin{enumerate}
    \item N.\ Alon, \textit{Combinatorial Nullstellensatz}, Combin.\ Probab.\ Comput.\ 8 (1999), 7--29.
    \item N.\ Alon, \textit{Eigenvalues and expanders}, Combinatorica 6 (1986), 83--96.
    \item N.\ Alon and M.\ Tarsi, \textit{Colorings and orientations of graphs}, Combinatorica 12 (1992), 125--134.
    \item N.\ Alon and J.\ Spencer, \textit{The Probabilistic Method}, Wiley, 4th ed., 2016.
    \item N.\ Alon, \textit{The algorithmic Lov\'asz Local Lemma}, Proc.\ SODA 2010.
\end{enumerate}

\subsection*{Expository Works}
\begin{enumerate}
    \item N.\ Alon, \textit{The combinatorial Nullstellensatz and its applications}, Proc.\ ICM 2002.
    \item L.\ Lov\'asz, \textit{Combinatorial Problems and Exercises}, AMS, 2nd ed., 2007.
    \item T.\ Tao and V.\ Vu, \textit{Additive Combinatorics}, Cambridge Univ.\ Press, 2006.
\end{enumerate}

